\documentclass[11pt]{article}
\usepackage[T1]{fontenc}
\usepackage{textcomp}
\usepackage[margin=0.75in]{geometry}
\usepackage{amsmath,amssymb,amsthm}
\usepackage{array,booktabs}
\usepackage{hyperref}
\setlength{\parindent}{0pt}
\newtheorem{theorem}{Theorem}
\theoremstyle{definition}
\newtheorem{definition}{Definition}
\title{Flow Proof Artifact: Bool-or-commutes}
\author{Generated by \texttt{flow doc proof}}
\date{Flow Proof Book}
\begin{document}
\maketitle
"A or B" gives the same result as "B or A".\\[0.5em]
\textbf{Source.} boole — https://en.wikipedia.org/wiki/Boolean\_algebra\\[0.5em]
\bigskip
\noindent{\large \textbf{Derived fact 1.} \textit{"A or B" gives the same result as "B or A"} \hfill boolean truth values \cdot disjunction \cdot order does not matter}\par
\smallskip\noindent\textit{Coordinate: boolean truth values \cdot disjunction \cdot order does not matter}\par
\smallskip\noindent\textit{Source: boole — https://en.wikipedia.org/wiki/Boolean\_algebra}\par
\medskip
\noindent\textbf{Goal.} "A or B" gives the same result as "B or A"\par
\noindent$\displaystyle \forall a \in \{\mathsf{true}, \mathsf{false}\} \forall b \in \{\mathsf{true}, \mathsf{false}\}\quad a \lor b = b \lor a$\par
\medskip
\noindent
\begin{tabular}{@{} >{\raggedright\arraybackslash}p{0.06\textwidth} >{\raggedright\arraybackslash}p{0.47\textwidth} >{\raggedleft\arraybackslash}p{0.41\textwidth} @{}}
\textbf{\#} & \textbf{Proof} & \textbf{Mathematics} \\
\midrule
\textbf{1.} & We prove that order does not matter for disjunction on boolean truth values. & \\[0.45em]
\textbf{2.} & We split into exhaustive cases — the claim must hold in each one. & \\[0.45em]
\textbf{3.} & Case 1 (see step 2): suppose a holds. & \\[0.45em]
\textbf{4.} & From step 3, this implies the disjunction of a and b equals the disjunction of b and a in this case. & $\displaystyle a \lor b = b \lor a$ \\[0.45em]
\textbf{5.} & Case 2 (see step 2): suppose b holds. & \\[0.45em]
\textbf{6.} & From step 5, this implies the disjunction of a and b equals the disjunction of b and a in this case. & $\displaystyle a \lor b = b \lor a$ \\[0.45em]
\textbf{7.} & Case 3 (see step 2): neither disjunct holds. & \\[0.45em]
\textbf{8.} & From step 7, this implies the disjunction of a and b equals the disjunction of b and a. Together with the other cases (step 3, step 5, and step 7), the goal is discharged. Hence proven. & $\displaystyle a \lor b = b \lor a$ \\[0.45em]
\bottomrule
\end{tabular}
\label{thm:Bool-disjunction-order_does_not_matter}
\par\medskip\hrule
\end{document}
